% PeerJ Computer Science review manuscript.
% Adapted from the PeerJ/Overleaf LaTeX submission template.
\documentclass[fleqn,12pt,lineno]{wlpeerj}

\usepackage[T1]{fontenc}
\usepackage[authoryear,round]{natbib}
\usepackage[hyphens]{url}
\usepackage{xurl}
\usepackage[hidelinks]{hyperref}
\usepackage[section]{placeins}
\usepackage{array}
\geometry{letterpaper,margin=1in}
\urlstyle{same}
\setlength{\emergencystretch}{1.8em}
\providecommand{\tightlist}{\setlength{\itemsep}{0pt}\setlength{\parskip}{0pt}}

\title{Three Protocol Choices Shape Measured Fidelity of AI Inpainting on Museum Paintings}
\author[1]{Wei Meng}
\author[2]{Wei Xiong}
\author[3]{Zhanhao Liang}
\affil[1]{Guilin Normal College, Guilin, Guangxi, China}
\affil[2]{Guilin Institute of Information Technology, Guilin, Guangxi, China}
\affil[3]{School of Art and Design, Guangzhou Polytechnic University, Guangzhou 511483, China}
\corrauthor[2]{Wei Xiong}{xiongwei@guit.edu.cn}

\begin{abstract}
\textbf{Background.} Full-reference benchmarks of off-the-shelf painting inpainting depend on three pipeline choices: latent mask downsampling, hole prefill color, and safety-filter denominator accounting.

\textbf{Methods.} We quantify these choices across four fixed models---Large Mask Inpainting (LaMa), Mask-Aware Transformer (MAT), PowerPaint v1, and Stable Diffusion XL (SDXL) Inpainting---on 120 open-access paintings from the Cleveland Museum of Art and Art Institute of Chicago spanning four medium-and-support groups and six procedural losses, using crops taken without further resampling from each museum\textquotesingle s published rendition.

\textbf{Results.} Nearest-neighbor latent downsampling misses a mean of 59.5\% of 1\%-crack pixels across 60 Cleveland works; grid alignment cuts PowerPaint v1 masked mean absolute error (MAE) by about 45\% across 40 complete paired works (safety checker on) while reducing availability. At SDXL\textquotesingle s model-card example strength, switching black prefill to a ring mean shifts overall center-masked MAE by 0.0725 and reverses luminance bias. PowerPaint\textquotesingle s safety checker withholds 21.9\% of Cleveland conditions versus 2.8\% in Chicago. Under the controlled protocol, LaMa \ensuremath{\approx} MAT \textless{} PowerPaint v1 \textless{} SDXL held across both museums, four media, and six losses without reversal.

\textbf{Conclusions.} In these cohorts, the best-performing model did not depend on painting medium; reporting these three choices is essential for reproducible evaluation.
\end{abstract}

\begin{document}
\maketitle
\raggedright
\linenumbers
\noindent\textbf{Keywords:} image inpainting; digital painting restoration; evaluation protocols; museum digitizations; full-reference metrics; latent diffusion

\section{Introduction}\label{introduction}

Digital image inpainting with artificial intelligence (AI) models offers a fast way to visualize completions for losses in digitized paintings, from fine craquelure and paint flaking to broad lacunae. In cultural heritage contexts, however, visual plausibility alone is insufficient. Conservation ethics treat loss compensation as a reversible, detectable, and documented intervention grounded in respect for original material \citep{icomcc2008, aic1994}, a principle recently demonstrated in physical practice through reversible digitally printed masks \citep{kachkine2025}. When digital inpainting supports documentation, visualization, or restoration planning, the central technical question is reconstruction fidelity: how closely a model recovers missing local color, brushwork, and structure rather than inventing plausible but unfaithful content. Because authentic historical losses lack surviving ground truth, controlled evaluation relies on masking intact museum digitizations and scoring predicted pixels directly against the recorded artwork using automatic full-reference distances.

Existing painting-inpainting studies have compared neural models under simulated or transferred losses, yet three low-level protocol choices in modern pipelines can substantially alter measured fidelity and the scored cohort (Fig.~\ref{figure:1}). First, the tested Stable Diffusion XL (SDXL) Inpainting and PowerPaint pipelines downsample pixel-space binary masks by a factor of eight using nearest-neighbor interpolation \citep{diffuserssdxl0352, powerpaintv1src}. Subsampling onto an 8-pixel grid can omit narrow structures from explicit mask conditioning, a broader issue in latent-space inpainting \citep{shamsolmoali2025tdpaint}; on 1\% painting cracks and flakes, it drops over half of the intended hole pixels. Second, when SDXL runs with denoising strength below 1.0, as in its model card\textquotesingle s 0.99 example \citep{sdxlinpaint}, the pipeline encodes the pre-masked input image into the initial noisy latent \citep{diffuserssdxl0352}. Because ground-truth hole pixels must be erased prior to encoding, the replacement color used to prefill the hole can affect generated outputs---a qualitative sensitivity discussed in Diffusers development \citep{diffusers8450}, whose quantitative luminance and error effects on paintings we measure here. Third, bundled safety checkers can withhold generated outputs; scoring only returned images compares models on different, pipeline-filtered subsets whose size varies sharply across collections. Alongside these protocol gaps, prior multi-model benchmarks either aggregate across heterogeneous artworks \citep{hu2025imagest} or stratify by visual subject such as portraits and landscapes \citep{controlled300}, leaving practitioners uncertain whether model rankings shift across distinct physical media and supports such as oil on canvas, oil on wood, watercolor on paper, and ink on paper.

\begin{figure}[!t]
\centering
\includegraphics[width=\linewidth]{figures/figure-1-teaser.pdf}
\caption{Three protocol choices change measured fidelity and sample coverage in artificial intelligence (AI) painting inpainting, while the four-model ordering holds across cohorts. \textbf{(a) Mask alignment:} default 8\ensuremath{\times} nearest-neighbor downsampling leaves part of a thin 1\% crack outside the model mask (red; 57.8\% missed in this single example from Edward Hopper\textquotesingle s \emph{Lighthouse Village}, main full-size crop; mean 59.5\% across 60 Cleveland Museum of Art (CMA) works); 8 \ensuremath{\times} 8 grid alignment reaches every hole pixel and lowers PowerPaint v1 masked mean absolute error (MAE) by about 45\% across 40 complete paired works (240 conditions, safety checker on). \textbf{(b) Hole prefill:} at Stable Diffusion XL (SDXL) \texttt{strength\ =\ 0.99}, black, mid-gray, and ring-mean prefill yield mean signed luminance biases of \ensuremath{-}0.085, +0.055, and +0.034 across 16 native-scale works (96 conditions), illustrated on a 5\% contiguous-loss case enlarged around the hole from the same Hopper artwork (prefill-arm museum rendition). \textbf{(c) Withheld outputs:} PowerPaint\textquotesingle s safety checker withholds outputs (red) in 79 of 360 CMA and 10 of 360 Art Institute of Chicago (AIC) conditions. \textbf{(d) Model ordering:} under the controlled protocol, Large Mask Inpainting (LaMa) \ensuremath{\approx} Mask-Aware Transformer (MAT) \textless{} PowerPaint v1 \textless{} SDXL holds across two museums, four media, and six loss conditions.}\label{figure:1}
\end{figure}

In this study, we quantify the effects of all three protocol choices and benchmark four fixed off-the-shelf inpainting models---two non-diffusion architectures, Large Mask Inpainting (LaMa) \citep{suvorov2021} and Mask-Aware Transformer (MAT) \citep{li2022mat}, and two latent diffusion models, PowerPaint v1 \citep{zhuang2023} and SDXL Inpainting \citep{sdxlinpaint}---under a pre-specified protocol, committed to version control before each inference run (Fig.~\ref{figure:2}). Across 120 open-access paintings equally balanced across four medium-and-support groups from the Cleveland Museum of Art (CMA, 60 works) and an independently sampled replication cohort from the Art Institute of Chicago (AIC, 60 works), we extract 512 \ensuremath{\times} 512 crops taken without further resampling from each museum\textquotesingle s published rendition and apply six shared procedural loss conditions spanning cracks, scattered flakes, and contiguous losses at 1\% and 5\% area fractions. We isolate latent-grid alignment through a full 60-work Cleveland paired comparison supported by our earlier 16-work pilot, isolate prefill color through a univariate three-color comparison (black, mid-gray, surrounding ring mean) plus a full-cohort SDXL sensitivity arm, and characterize safety-filter attrition through planned-denominator tracking and a conditional replay of all withheld Cleveland cases.

Across two museum cohorts and targeted diagnostic, prefill, and replay experiments, our study makes four main contributions:

\begin{enumerate}
\def\labelenumi{\arabic{enumi}.}
\tightlist
\item
  \textbf{Quantification of latent-space mask under-coverage and alignment tradeoffs.} Across the Cleveland cohort, default nearest-neighbor latent downsampling omits a mean of 59.5\% of narrow crack pixels; aligning masks to the latent grid lowers PowerPaint v1 masked mean absolute error (MAE) by about 45\% on complete paired works while reducing its scorable output rate, whereas SDXL shows mixed condition-dependent shifts.
\item
  \textbf{Quantification of prefill-color sensitivity under partial-strength denoising.} Fixing SDXL at the model card\textquotesingle s 0.99 example settings and switching hole prefill from black to a surrounding ring mean shifts overall center-masked MAE by 0.0725 (with larger condition-mean shifts on broad losses) and reverses luminance bias; pairing surrounding-ring prefill with those example settings substantially lowers Cleveland SDXL MAE while leaving a wide gap above LaMa.
\item
  \textbf{Quantification of cohort-dependent safety-filter withholding.} PowerPaint v1 withholds 21.9\% of planned conditions in Cleveland versus 2.8\% in Chicago under identical settings, and replayed Cleveland conditions exhibit higher mean error than complete safety-on works, demonstrating why coverage denominators must accompany fidelity scores.
\item
  \textbf{Cross-museum, cross-medium stability of model rankings.} Under the controlled protocol, the masked-MAE ordering LaMa \(\approx\) MAT \(\!<\!\) PowerPaint v1 \(\!<\!\) SDXL holds identically across both museums, all four medium-and-support strata, and all six damage conditions (with overall LaMa--MAT differences inside our 0.01 MAE practical margin); the three protocol choices changed error magnitudes and the scored sample without reversing this ordering in our tests, whereas between-group gap magnitudes varied across collections.
\end{enumerate}

\section{Related Work}\label{related-work}

\subsection{General-purpose image inpainting pipelines}\label{general-purpose-image-inpainting-pipelines}

Classical inpainting formulated hole filling as local continuation from known boundaries, propagating isophotes via partial differential equations \citep{bertalmio2000} or copying best-matching patches by structure-prioritized exemplar synthesis \citep{criminisi2004}. Deep learning shifted inpainting toward learned global priors, beginning with context encoders combining reconstruction and adversarial losses \citep{pathak2016} and extending to arbitrary hole geometries via partial convolutions \citep{liu2018} and gated convolutions \citep{yu2018}. To bridge wide gaps at higher resolutions, subsequent non-diffusion architectures introduced co-modulated generative adversarial networks \citep{zhao2021}, fast Fourier convolutions with image-wide receptive fields in LaMa \citep{suvorov2021}, and mask-aware transformer attention in MAT \citep{li2022mat}. Because LaMa and MAT operate directly on pixel-grid inputs without text conditioning or content filters, they provide strong non-diffusion baselines for local continuation and run deterministically under our fixed-seed configuration.

A second paradigm performs conditional generation inside the compressed latent space of an autoencoder \citep{rombach2021} or via iterative diffusion resampling \citep{lugmayr2022}. Widely used open checkpoints include SDXL Inpainting 0.1, trained at 1024 \ensuremath{\times} 1024 resolution with mask and masked-image UNet conditioning channels \citep{sdxlinpaint}; PowerPaint v1, which adds learnable task prompts to separate context-aware filling from object insertion and removal \citep{zhuang2023}; dual-branch architectures such as BrushNet \citep{ju2024}; and large rectified-flow transformers such as FLUX.1 Fill \citep{fluxfill}. In the latent diffusion pipelines tested here, three operational and deployment choices absent from our pixel-space baselines govern what is scored: spatial quantization of fine masks onto an 8\ensuremath{\times} downsampled latent grid \citep{shamsolmoali2025tdpaint, ju2024, diffuserssdxl0352, powerpaintv1src}, coupling between pre-masked red--green--blue (RGB) inputs and initial noisy latents whenever \texttt{strength\ \textless{}\ 1.0} \citep{diffuserssdxl0352, diffusers8450}, and an optional safety filter enabled by default in the tested PowerPaint deployment \citep{zhuang2023, powerpaintv1src}. Our study quantifies how these three pipeline choices affect full-reference fidelity and sample availability on digitized paintings.

\subsection{Digital painting restoration, damage modeling, and benchmarks}\label{digital-painting-restoration-damage-modeling-and-benchmarks}

Early computational restoration studies demonstrated virtual crack detection and inpainting on masterworks such as the Ghent Altarpiece while noting the lack of surviving ground truth for authentic damage \citep{cornelis2013}. Subsequent work paired domain-specific models with heritage datasets, including Dunhuang mural benchmarks using random-walk or real-damage masks \citep{yu2019dunhuang, shao2023, xu2024dunhuang}, deep image prior formulations for frescoes and style-coherent filling \citep{merizzi2024, morotti2024}, dual-encoder networks for ancient Chinese paintings \citep{sun2024}, super-resolution and textile restoration pipelines \citep{hu2024, wang2024silk}, fine-tuned diffusion for frescoes \citep{smolka2025}, and synthetic craquelure generation \citep{cuchguillen2026}. In parallel, pixel-level damage segmentation datasets such as ARTeFACT catalog real analogue damage across heritage media \citep{ivanova2024artefact}, while several restoration papers transplant real damage silhouettes onto intact Chinese paintings \citep{zhang2026inkrenew, xue2025}. Creative painterly editing frameworks that insert new objects or text-guided content pursue a different goal from conservation-oriented reconstruction \citep{xu2023refpaint, fanelli2024}.

Four recent benchmarks are closest to our multi-model evaluation. Cipolina-Kun et al. compared CoModGAN, LaMa, and GLIDE on transformations of a single Escher lithograph plus a qualitative fresco example \citep{cipolinakun2022}, leaving multi-artwork and cross-medium generalization untested. Imagest benchmarked a style-conditioned generator against VQGAN, Stable Diffusion 1.4/1.5, and DALL·E 2 on 100 artificially damaged paintings cropped or resized to 256 \ensuremath{\times} 256 pixels, reporting aggregate full-reference metrics without medium or support stratification \citep{hu2025imagest}. MDTD-ArtIR evaluated blind restoration and editing models on 1,000 WikiArt paintings degraded by semi-transparent texture overlays rather than binary missing-region masks \citep{vijendran2026}. Most closely, Controlled-300---an unreviewed Zenodo master\textquotesingle s research archive \citep{controlled300}---compared Telea, LaMa, HINT, and Stable Diffusion across 300 paintings grouped by visual content (portraits, landscapes, architecture, abstract, and high texture) alongside a 35-work fixed-area mask robustness check, reporting low masked MAE for LaMa. Relative to these precedents, we quantify latent-grid mask coverage, prefill initialization, and safety-filter denominators, stratify explicitly by physical medium and support rather than visual genre, and test replication across two independent museum open-access collections without downsampling museum renditions into uniform 256-pixel thumbnails.

\subsection{Full-reference evaluation metrics and conservation ethics}\label{full-reference-evaluation-metrics-and-conservation-ethics}

Automatic restoration evaluation must distinguish pixel-level distortion from perceptual or structural similarity, a tension formalized at the distribution level by the perception--distortion tradeoff \citep{blau2018}. Pixelwise RGB error (MAE) directly measures local numeric agreement, while CIEDE2000 (\texttt{\ensuremath{\Delta E_{00}}}) provides an unsigned, perceptually weighted color distance in CIE 1976 L\emph{a}b* (CIELAB) space \citep{luo2001ciede2000, sharma2005ciede2000}. Learned full-reference metrics extend comparison to mid- and high-level representations: Learned Perceptual Image Patch Similarity (LPIPS) aligns with human pairwise similarity judgments \citep{zhang2018lpips}, whereas Deep Image Structure and Texture Similarity (DISTS) unifies structure and texture similarity with partial tolerance to stochastic texture resampling \citep{ding2020dists}. More recent work targets localized failure modes more specifically, including MISATO\textquotesingle s \texttt{CLIP@mask} for masked semantic alignment and \texttt{Gradient@edge} for boundary color seams \citep{wang2025misato}, high-level fidelity proxies in super-resolution \citep{rocafort2025}, controllable hallucination testing in HalluGen \citep{kim2026hallugen}, and discussions of iconographic hallucination in mural restoration \citep{ren2026mapgr}.

Conservation charters establish why transparent, auditable protocols matter even in purely digital evaluations. The 2008 International Council of Museums Committee for Conservation (ICOM-CC) terminology resolution defines restoration as actions undertaken with respect for original material \citep{icomcc2008}, and Section 23 of the American Institute for Conservation Guidelines for Practice requires compensation for loss to be documented, detectable by common examination methods, and reversible \citep{aic1994}. Translating these norms to a computational benchmark motivates preserving exact unmasked original pixels outside the scoring hole, logging immutable inputs and outputs, and reporting all safety-filtered cases alongside full-reference scores rather than silently discarding failures.

\section{Materials and Methods}\label{materials-and-methods}

\subsection{Study overview and pre-specified workflow}\label{study-overview-and-pre-specified-workflow}

To separate architectural differences from evaluation artifacts, we built a modular, pre-specified benchmark pipeline (Fig.~\ref{figure:2}). Every experimental stage---the center-masked pilot, the 16-work native-scale pilot and latent-alignment correction, the SDXL parameter and univariate prefill studies, the 60-work Cleveland main cohort, the Cleveland review and conditional replay, the 60-work Chicago replication, and the full 60-work Cleveland original-mask comparison---was specified in a version-controlled protocol committed before graphics processing unit (GPU) inference, with run windows and artifact hashes logged for verification. Beginning with the full-cohort Cleveland original-mask comparison, each run sidecar additionally recorded the executed Git commit hash, working-tree cleanliness, binary diff digest, and Python source hash at launch. All results in this study characterize reconstruction fidelity and output availability for the pinned off-the-shelf checkpoints, prompts, single-seed main settings, and 512-pixel procedural loss protocols described below. Full implementation, sampling audit, and checksum details are provided in Supplementary Materials (\nolinkurl{docs/paper/SUPPLEMENT.md}).

\begin{figure}[!htb]
\centering
\includegraphics[width=.63\linewidth]{figures/figure-2-pipeline.pdf}
\caption{Controlled two-museum evaluation pipeline. Crops from two open-access collections (60 works each, 15 per medium-and-support group) receive six shared procedural losses per artwork. All four pinned models share the 8 \ensuremath{\times} 8 grid-aligned input aperture (a); the SDXL sensitivity varies hole prefill color (b); generated pixels are kept only inside the original hole for full-reference scoring, while planned, scored, and withheld conditions are counted against 360 planned conditions per model per museum (c). Colored elements highlight the three protocol controls; all protocols are fixed prior to inference.}\label{figure:2}
\end{figure}

\subsection{Museum cohorts, sampling, and crop preparation}\label{museum-cohorts-sampling-and-crop-preparation}

We assembled two independent open-access painting cohorts of 60 artworks each (120 total), evenly balanced across four physical medium-and-support groups (15 works per group): oil on canvas, oil on wood support, watercolor on paper, and ink on paper.

\begin{itemize}
\tightlist
\item
  \textbf{Cleveland Museum of Art (CMA) cohort.} Sixty Creative Commons Zero (CC0) paintings were sampled deterministically from Cleveland\textquotesingle s open-access collection using the museum\textquotesingle s full-size downloadable image renditions \citetext{\citeauthor{cma-openaccess}, n.d.}. Sixteen works overlapped with our earlier native-scale pilot. One frozen crop (\emph{The Trinity}) was uniformly black RGB \texttt{(0,\ 0,\ 0)} in the museum source file; under our pre-specified rule against post-selection replacement, it was retained in the main cohort and verified in an exploratory exclusion check.
\item
  \textbf{Art Institute of Chicago (AIC) independent replication cohort.} Sixty public-domain paintings were sampled independently from a frozen AIC catalog export \citetext{\citeauthor{aic-openaccess}, n.d.} using deterministic Secure Hash Algorithm 256 (SHA-256) ordering within each medium-and-support stratum, restricting selection to verified official 1686-pixel-wide International Image Interoperability Framework (IIIF) Joint Photographic Experts Group (JPEG) renditions, at most one work per identifiable primary artist per group, and zero image or object overlap with Cleveland. One watercolor (\texttt{aic-114337}) painted on paper mounted onto card met the frozen catalog filter and was verified via a separate exclusion check (Supplementary Note S1 details candidate screening and exclusions).
\end{itemize}

After Exchangeable Image File Format (EXIF) orientation correction and conversion of embedded color profiles to standard RGB (sRGB), each artwork contributed one 8-bit RGB 512 \ensuremath{\times} 512 crop at a fixed pseudorandom integer coordinate, \textbf{taken without further resampling from each museum\textquotesingle s published rendition}. Each crop received six procedural binary masks crossing three damage morphologies---a thin tortuous connected crack, dispersed small paint flakes, and a single contiguous irregular loss---with two target area fractions (1\% and 5\% of the 512 \ensuremath{\times} 512 crop), totaling 360 planned conditions per museum. The six masks assigned to within-group ranks 0--14 were identical across all four medium-and-support groups and identical between Cleveland and Chicago, ensuring exact geometric matching across strata and collections.

\subsection{Pinned models and the three controlled protocol experiments}\label{pinned-models-and-the-three-controlled-protocol-experiments}

We evaluated four publicly released inpainting models on a single NVIDIA RTX 3090 (24 GiB video random-access memory, VRAM) using fixed seed \texttt{2026} in all main runs (checkpoint versions, adapter notes, and FLUX.1 Fill \citep{fluxfill} pre-screening are detailed in Supplementary Note S2):

\begin{itemize}
\tightlist
\item
  \textbf{LaMa} \citep{suvorov2021}: IOPaint TorchScript release, standard 512-pixel inference without iterative refinement.
\item
  \textbf{MAT} \citep{li2022mat}: IOPaint Places-512 release (\texttt{truncation\_psi\ =\ 1}) with deterministic per-condition latent seeding.
\item
  \textbf{PowerPaint v1} \citep{zhuang2023}: pinned author release, Denoising Diffusion Implicit Models (DDIM) scheduler, 512 resolution, 25 steps, classifier-free guidance (CFG) 7.5, \texttt{strength\ =\ 1.0}, positive context task token (\texttt{P\_ctxt}) paired with a single generic painting-restoration prompt, negative object token (\texttt{P\_obj}), and the built-in safety checker kept enabled in all main arms.
\item
  \textbf{SDXL Inpainting 0.1} \citep{sdxlinpaint}: pinned Diffusers checkpoint, 9-channel UNet, EulerDiscreteScheduler, 512 resolution, the same generic painting-restoration prompt, 25 requested steps, CFG 7.5, and \texttt{strength\ =\ 1.0} in main arms.
\end{itemize}

\textbf{Protocol Control 1: Latent-space mask alignment and shared input aperture.} Both SDXL Inpainting and PowerPaint downsample binary masks by 8\ensuremath{\times} via nearest-neighbor interpolation onto a 64 \ensuremath{\times} 64 latent grid at 512-pixel inference \citep{diffuserssdxl0352, powerpaintv1src}. In our main comparison, we enforced a shared latent-aligned aperture across all four models: every 8 \ensuremath{\times} 8 pixel block intersecting the original scoring mask was marked as masked in the model input, and pixels inside this expanded mask were zeroed (black) before inference. Final compositing placed generated pixels strictly inside the original unexpanded scoring mask and preserved exact original pixels everywhere outside (verified pixel-for-pixel across all outputs). To measure the effect of mask alignment directly at full scale, we ran a pre-specified paired comparison on all 60 Cleveland works (360 conditions each for PowerPaint v1 and SDXL) feeding the unexpanded original mask directly into both diffusion pipelines while keeping all other settings fixed. We also ran original-mask input sensitivity arms for LaMa and MAT in both Cleveland and Chicago (720 runs per museum) and retained our earlier 16-work native-scale and 60-work center-masked alignment runs as supporting evidence.

\textbf{Protocol Control 2: Prefill initialization and SDXL parameter sensitivity.} In Diffusers\textquotesingle{} SDXL Inpainting pipeline \citep{diffuserssdxl0352}, \texttt{strength\ =\ 1.0} initializes denoising from pure Gaussian noise while keeping zero-centered masked-image conditioning. When \texttt{strength\ \textless{}\ 1.0}---such as the model card\textquotesingle s example settings of \texttt{strength\ =\ 0.99}, CFG 8.0, and 20 requested steps \citep{sdxlinpaint}, for which the pinned Diffusers v0.35.2 pipeline executes 19 effective denoising steps \citep{diffuserssdxl0352}---the pipeline additionally encodes the full pre-masked RGB input through the variational autoencoder (VAE) and adds noise to that latent (Supplementary Note S3 gives the Euler initialization scaling). Because ground-truth pixels must be erased before VAE encoding, the hole prefill color can directly shift generated luminance and error.

To isolate prefill color as a single variable, we fixed SDXL at 512 resolution, \texttt{strength\ =\ 0.99}, CFG 8.0, 20 requested steps, and seed \texttt{2026}, varying only the RGB color inside the expanded hole across three arms: black (RGB 0), mid-gray (RGB 128), and an 8-pixel surrounding ring mean (the single per-image rounded sRGB mean of visible pixels within an 8-pixel neighborhood strictly outside the expanded mask). This univariate test ran on both the 60 center-masked works (5\% irregular central hole) and the 16 native-scale pilot works (96 conditions per color). We then evaluated the strongest configuration (\texttt{0.99\ /\ 8.0\ /\ 20} + ring-color prefill) across all 360 Cleveland main conditions to test whether favorable SDXL initialization closes the gap to non-diffusion baselines. Additional single-seed sensitivities on the 60 center-masked works tested 1024-pixel inference and an empty prompt. We also visually audited 12 randomly sampled Cleveland SDXL main cases via three-column contact sheets (original, LaMa, SDXL) and reproduced two cases pixel-for-pixel to verify determinism.

\textbf{Protocol Control 3: Safety-filter denominators and conditional replay.} Whenever PowerPaint\textquotesingle s built-in safety checker withheld an output, the run was repeated once with the identical seed and inputs to check for transient runtime failures; persistent flags were recorded as withheld cases rather than scored as black images, and kept in the planned denominator. To assess how Cleveland\textquotesingle s 79 withheld conditions (19 works) affected the four-model comparison, we audited museum catalog subject descriptions of those 19 paintings and executed a single-pass conditional replay of exactly those 79 conditions with the safety checker disabled (identical crops, masks, seed, and settings), evaluating both the 79-condition replayed subset and the combined 60-work / 360-condition mixed cohort alongside the primary safety-on cohort.

\subsection{Metrics and statistical analysis}\label{metrics-and-statistical-analysis}

Our primary fidelity metric is \textbf{masked MAE}: the mean absolute RGB difference between the reconstructed and original crops computed strictly over pixels inside the original unexpanded mask, divided by 255 to lie on \texttt{{[}0,\ 1{]}}. Four secondary full-reference metrics capture chromatic accuracy, perceptual similarity, and texture/structure preservation without human raters:

\begin{enumerate}
\def\labelenumi{\arabic{enumi}.}
\tightlist
\item
  \textbf{Pixel-pair mean CIEDE2000 (\texttt{\ensuremath{\Delta E_{00}}})} inside the original mask (unsigned perceptual color distance; D65 illuminant, CIE 1931 2° observer) \citep{luo2001ciede2000, sharma2005ciede2000}.
\item
  \textbf{CIEDE2000 of masked mean CIELAB colors}, an unsigned distance measuring the magnitude of net color difference between the predicted and reference hole averages without shift direction.
\item
  \textbf{Local LPIPS} (AlexNet v0.1) \citep{zhang2018lpips} on the original mask\textquotesingle s tight bounding box padded by an 8-pixel margin (clipped to the 512 \ensuremath{\times} 512 boundary).
\item
  \textbf{Local DISTS} \citep{ding2020dists} on the same padded bounding box.
\end{enumerate}

For prefill analyses, we additionally computed \textbf{signed luminance bias}: the mean difference (output minus reference) in Rec. 709 weighted sRGB luminance \texttt{{[}0,\ 1{]}} inside the original mask, where negative values denote darkening and positive values denote brightening. For all distance metrics, lower values indicate closer agreement with the reference crop.

Within each museum, overall artwork-level scores weighted the six damage conditions equally and included only artworks with all six conditions available for the comparison at hand. Four-model summary tables therefore use the common complete-work denominator (41/60 in Cleveland, 54/60 in Chicago), whereas pairwise comparisons against LaMa retain every artwork complete for that specific model pair (60/60 for MAT and SDXL in both museums; 41 and 54 for PowerPaint). Per-condition summaries use all artworks complete across all models in that comparison. Uncertainty was quantified via 10,000 artwork-level cluster bootstrap resamples, yielding unadjusted 95\% percentile confidence intervals (CIs). Before running the main experiments, we pre-specified practical equivalence margins of \textbf{0.01 masked MAE} (about 2.55 gray levels per 8-bit channel) and \textbf{0.02 local LPIPS}; medium-and-support contrasts and per-condition breakdowns are interpreted descriptively.

\subsection{Use of Generative AI Tools}\label{use-of-generative-ai-tools}

Claude (Anthropic) assisted with protocol review and figure generation. Gemini (Google) assisted with English drafting and editing. Codex (OpenAI) assisted with experimental and figure-generation code, literature verification, and statistical verification scripts \emph{(author to confirm specific tool version details and exact scope of personal review prior to submission)}. These uses supported implementation, auditability, and clear reporting of the evaluation. Figures~\ref{figure:1} and~\ref{figure:3} show saved inpainting outputs from the evaluated models under the fixed experimental protocol. The research questions, methodological decisions, pre-inference protocol approvals, and final scientific interpretations were decided by the research team. No generative AI system was used to fabricate measurements or modify raw model outputs; all reported metrics, denominators, and confidence intervals were recomputed programmatically by deterministic scripts directly from frozen run outputs and verified against source manifests.

\section{Results}\label{results}

\subsection{Model rankings replicate identically across both museum collections}\label{model-rankings-replicate-identically-across-both-museum-collections}

Under the controlled protocol, the four-model ordering was identical in both Cleveland and Chicago for our two thresholded metrics, masked MAE and local LPIPS: \textbf{LaMa \(\approx\) MAT \(\!<\!\) PowerPaint v1 \(\!<\!\) SDXL} (Table~\ref{table:1}). Point-estimate means followed the same sequence (LaMa, MAT, PowerPaint v1, SDXL) for pixel-pair \texttt{\ensuremath{\Delta E_{00}}} and local DISTS, while MAT scored slightly lower than LaMa on unsigned mean-Lab \texttt{\ensuremath{\Delta E_{00}}} in both museums. On complete four-model cohorts, LaMa achieved a masked MAE of 0.0315 in Cleveland (41 complete works, 246 conditions) and 0.0387 in Chicago (54 complete works, 324 conditions), closely matched by MAT at 0.0339 and 0.0421. PowerPaint v1 increased masked MAE (0.0581 and 0.0650) and more than doubled local LPIPS and DISTS relative to LaMa, while SDXL under the main \texttt{strength\ =\ 1.0} configuration produced substantially larger errors across all five metrics (MAE 0.3076 in Cleveland, 0.2709 in Chicago; mean-Lab \texttt{\ensuremath{\Delta E_{00}}} exceeding 26 in both collections). Figure~\ref{figure:3} shows representative reconstructions across crack, flake, and contiguous-loss conditions.

\begin{table}[!htb]
\caption{Mean full-reference reconstruction distances across both museum cohorts. Results are evaluated under the shared 8 \ensuremath{\times} 8 expanded input aperture and scored strictly inside the original mask (or its 8-pixel padded bounding box for perceptual metrics) for Large Mask Inpainting (LaMa), Mask-Aware Transformer (MAT), PowerPaint v1, and Stable Diffusion XL (SDXL) Inpainting. Metrics report masked mean absolute error (MAE), local Learned Perceptual Image Patch Similarity (LPIPS), pixel-pair mean and masked mean-Lab CIEDE2000 (\ensuremath{\Delta E_{00}}) in CIE 1976 L\emph{a}b* (CIELAB) space, and local Deep Image Structure and Texture Similarity (DISTS). Complete four-model denominators comprise \texttt{n\ =\ 41} works (246 conditions) in Cleveland and \texttt{n\ =\ 54} works (324 conditions) in Chicago; down arrows (\ensuremath{\downarrow}) indicate that lower values reflect closer fidelity.}\label{table:1}
\centering
\fontsize{10}{12}\selectfont
\setlength{\tabcolsep}{3pt}
\renewcommand{\arraystretch}{1.2}
\begin{tabular}{>{\raggedright\arraybackslash}p{\dimexpr0.170\linewidth-2.380\tabcolsep\relax}>{\raggedright\arraybackslash}p{\dimexpr0.155\linewidth-2.170\tabcolsep\relax}>{\raggedright\arraybackslash}p{\dimexpr0.112\linewidth-1.568\tabcolsep\relax}>{\raggedright\arraybackslash}p{\dimexpr0.112\linewidth-1.568\tabcolsep\relax}>{\raggedright\arraybackslash}p{\dimexpr0.155\linewidth-2.170\tabcolsep\relax}>{\raggedright\arraybackslash}p{\dimexpr0.166\linewidth-2.324\tabcolsep\relax}>{\raggedright\arraybackslash}p{\dimexpr0.130\linewidth-1.820\tabcolsep\relax}}
\toprule
\textbf{Collection (complete works)} & \textbf{Model} & \textbf{Masked MAE \ensuremath{\downarrow}} & \textbf{Local LPIPS \ensuremath{\downarrow}} & \textbf{Pixel-pair mean \ensuremath{\Delta E_{00}} \ensuremath{\downarrow}} & \textbf{Mean-Lab \ensuremath{\Delta E_{00}} \ensuremath{\downarrow}} & \textbf{Local DISTS \ensuremath{\downarrow}} \\
\midrule
Cleveland (n = 41) & LaMa & 0.0315 & 0.0643 & 3.7305 & 0.7947 & 0.0491 \\
Cleveland (n = 41) & MAT & 0.0339 & 0.0761 & 3.9414 & 0.7154 & 0.0542 \\
Cleveland (n = 41) & PowerPaint v1 & 0.0581 & 0.1696 & 6.2781 & 3.5030 & 0.1417 \\
Cleveland (n = 41) & SDXL (main) & 0.3076 & 0.4930 & 31.9019 & 30.5106 & 0.3391 \\
Chicago (n = 54) & LaMa & 0.0387 & 0.0685 & 4.1787 & 0.9272 & 0.0570 \\
Chicago (n = 54) & MAT & 0.0421 & 0.0773 & 4.4747 & 0.9173 & 0.0595 \\
Chicago (n = 54) & PowerPaint v1 & 0.0650 & 0.1482 & 6.9000 & 3.1668 & 0.1234 \\
Chicago (n = 54) & SDXL (main) & 0.2709 & 0.4181 & 28.2618 & 26.4494 & 0.2906 \\
\bottomrule
\end{tabular}
\end{table}


Artwork-level paired comparisons against LaMa confirm both the practical equivalence of the two non-diffusion baselines and the size of the diffusion gaps (Table~\ref{table:2}). Across all 60 works in each museum, MAT differed from LaMa by +0.0022 {[}0.0015, 0.0030{]} MAE in Cleveland and +0.0032 {[}0.0023, 0.0041{]} in Chicago; the corresponding local LPIPS differences (+0.0102 {[}0.0070, 0.0135{]} and +0.0088 {[}0.0066, 0.0110{]}) also lay entirely inside our pre-specified 0.01 MAE and 0.02 LPIPS practical margins. By contrast, PowerPaint v1 exceeded LaMa by +0.0266 {[}0.0203, 0.0337{]} MAE in Cleveland and +0.0263 {[}0.0221, 0.0310{]} in Chicago, and SDXL exceeded LaMa by +0.2690 {[}0.2203, 0.3186{]} and +0.2480 {[}0.2035, 0.2956{]}.

\begin{table}[!htb]
\caption{Artwork-level paired reconstruction differences relative to LaMa. Values report candidate minus LaMa differences with unadjusted 95\% bootstrap percentile confidence intervals (CI; 10,000 artwork resamples); positive values indicate higher error than LaMa.}\label{table:2}
\centering
\fontsize{10}{12}\selectfont
\setlength{\tabcolsep}{3pt}
\renewcommand{\arraystretch}{1.2}
\begin{tabular}{>{\raggedright\arraybackslash}p{\dimexpr0.110\linewidth-1.100\tabcolsep\relax}>{\raggedright\arraybackslash}p{\dimexpr0.170\linewidth-1.700\tabcolsep\relax}>{\raggedright\arraybackslash}p{\dimexpr0.110\linewidth-1.100\tabcolsep\relax}>{\raggedright\arraybackslash}p{\dimexpr0.305\linewidth-3.050\tabcolsep\relax}>{\raggedright\arraybackslash}p{\dimexpr0.305\linewidth-3.050\tabcolsep\relax}}
\toprule
\textbf{Collection} & \textbf{Candidate model} & \textbf{Complete works} & \textbf{Masked MAE difference [95\% CI]} & \textbf{Local LPIPS difference [95\% CI]} \\
\midrule
Cleveland & MAT & 60 & +0.0022 [0.0015, 0.0030] & +0.0102 [0.0070, 0.0135] \\
Cleveland & PowerPaint v1 & 41 & +0.0266 [0.0203, 0.0337] & +0.1053 [0.0800, 0.1333] \\
Cleveland & SDXL (main) & 60 & +0.2690 [0.2203, 0.3186] & +0.4258 [0.3893, 0.4623] \\
Chicago & MAT & 60 & +0.0032 [0.0023, 0.0041] & +0.0088 [0.0066, 0.0110] \\
Chicago & PowerPaint v1 & 54 & +0.0263 [0.0221, 0.0310] & +0.0797 [0.0625, 0.0990] \\
Chicago & SDXL (main) & 60 & +0.2480 [0.2035, 0.2956] & +0.3695 [0.3304, 0.4104] \\
\bottomrule
\end{tabular}
\end{table}


The planned original-mask sensitivity for pixel-space models and exploratory edge-case exclusions preserved these findings. Giving LaMa and MAT the unexpanded original mask reduced masked MAE by 0.0040 {[}0.0033, 0.0047{]} and 0.0055 {[}0.0046, 0.0064{]} across all 60 Cleveland works (complete-cohort means 0.0276 and 0.0285 on 41 works), and by 0.0058 {[}0.0047, 0.0070{]} and 0.0081 {[}0.0066, 0.0097{]} across all 60 Chicago works (0.0325 and 0.0334 on 54 works), maintaining LaMa \(\approx\) MAT \(\!<\!\) PowerPaint v1 \(\!<\!\) SDXL. Excluding Cleveland\textquotesingle s uniformly black crop (\emph{The Trinity}, n = 59) yielded SDXL minus LaMa MAE differences of +0.2597 {[}0.2137, 0.3065{]} (main) and +0.0776 {[}0.0704, 0.0852{]} (ring-prefill sensitivity); excluding Chicago\textquotesingle s card-mounted watercolor (\texttt{aic-114337}) yielded MAE differences vs. LaMa of +0.0031 (MAT, n = 59), +0.0265 (PowerPaint, n = 53), and +0.2480 (SDXL, n = 59).

\begin{figure}[!htb]
\centering
\includegraphics[width=\linewidth]{figures/figure-3-examples.pdf}
\caption{Qualitative comparison across three damage geometries and four inpainting models. Deterministically selected 5\% crack, scattered-flake, and contiguous-loss examples from three distinct CC0 Cleveland artworks are shown across columns (reference crop, binary scoring mask, LaMa, MAT, PowerPaint v1, and SDXL Inpainting under main settings), with magnified detail boxes and per-tile masked MAE; generated pixels are composited strictly inside the original mask.}\label{figure:3}
\end{figure}

\subsection{Latent-space mask under-coverage drops over half of thin losses and alters both error and availability}\label{latent-space-mask-under-coverage-drops-over-half-of-thin-losses-and-alters-both-error-and-availability}

Across all 60 Cleveland works (360 planned conditions), default 8\ensuremath{\times} nearest-neighbor downsampling omitted \textbf{59.54\%} of scoring-mask pixels on 1\% cracks and \textbf{52.93\%} on 1\% flakes, compared with 22.46\% on 5\% cracks, 25.01\% on 5\% flakes, 13.85\% on 1\% contiguous losses, and 6.63\% on 5\% contiguous losses (Table~\ref{table:3}; Fig.~\ref{figure:4}). Expanding the input mask to cover every intersecting 8 \ensuremath{\times} 8 block reduced missed pixels to zero across all six geometries while increasing the masked input area to 2.649\ensuremath{\times} the original area on 1\% cracks and 2.870\ensuremath{\times} on 1\% flakes (versus 1.164\ensuremath{\times} on 5\% contiguous losses).

\begin{table}[!htb]
\caption{Latent-space mask under-coverage and paired alignment effects across the full 60-work Cleveland cohort (Protocol Factor 1). Entries report pixel under-coverage under default 8\ensuremath{\times} nearest-neighbor downsampling, stride-8 expanded-to-original input area ratios, and paired masked MAE changes (expanded minus original input, seed 2026) with 95\% bootstrap CIs. Overall rows use artworks complete across all six conditions under both apertures (\texttt{n\ =\ 40} works, 240 conditions for PowerPaint v1; \texttt{n\ =\ 60} works, 360 conditions for SDXL); em dashes (---) indicate geometry-specific ratios not aggregated into a single overall figure.}\label{table:3}
\centering
\fontsize{10}{12}\selectfont
\setlength{\tabcolsep}{3pt}
\renewcommand{\arraystretch}{1.2}
\begin{tabular}{>{\raggedright\arraybackslash}p{\dimexpr0.135\linewidth-1.350\tabcolsep\relax}>{\raggedright\arraybackslash}p{\dimexpr0.140\linewidth-1.400\tabcolsep\relax}>{\raggedright\arraybackslash}p{\dimexpr0.135\linewidth-1.350\tabcolsep\relax}>{\raggedright\arraybackslash}p{\dimexpr0.295\linewidth-2.950\tabcolsep\relax}>{\raggedright\arraybackslash}p{\dimexpr0.295\linewidth-2.950\tabcolsep\relax}}
\toprule
\textbf{Procedural condition} & \textbf{Missed hole pixels (original mask)} & \textbf{Stride-8 expanded / original area ratio} & \textbf{PowerPaint v1 MAE change [95\% CI] (n)} & \textbf{SDXL MAE change [95\% CI] (n = 60)} \\
\midrule
Crack 1\% & 59.54\% & 2.649 & \ensuremath{-}0.1494 [\ensuremath{-}0.1853, \ensuremath{-}0.1143] (n = 45) & \ensuremath{-}0.0704 [\ensuremath{-}0.0878, \ensuremath{-}0.0540] \\
Crack 5\% & 22.46\% & 1.566 & +0.0164 [0.0055, 0.0297] (n = 46) & \ensuremath{-}0.0034 [\ensuremath{-}0.0257, 0.0201] \\
Flake 1\% & 52.93\% & 2.870 & \ensuremath{-}0.1707 [\ensuremath{-}0.2015, \ensuremath{-}0.1409] (n = 47) & \ensuremath{-}0.0652 [\ensuremath{-}0.0811, \ensuremath{-}0.0507] \\
Flake 5\% & 25.01\% & 1.725 & +0.0082 [\ensuremath{-}0.0003, 0.0175] (n = 45) & +0.0594 [0.0459, 0.0737] \\
Contiguous loss 1\% & 13.85\% & 1.371 & \ensuremath{-}0.0016 [\ensuremath{-}0.0048, 0.0012] (n = 46) & +0.0269 [0.0024, 0.0533] \\
Contiguous loss 5\% & 6.63\% & 1.164 & +0.0206 [0.0067, 0.0385] (n = 46) & +0.0574 [0.0388, 0.0770] \\
\textbf{Overall paired works} & \textbf{—} & \textbf{—} & \textbf{\ensuremath{-}0.0471 [\ensuremath{-}0.0589, \ensuremath{-}0.0360] (0.1055 \ensuremath{\to} 0.0583; n = 40)} & \textbf{+0.0008 [\ensuremath{-}0.0081, 0.0101] (0.3007 \ensuremath{\to} 0.3015; n = 60)} \\
\bottomrule
\end{tabular}
\end{table}


In the full-cohort paired comparison, grid alignment produced a major fidelity improvement for PowerPaint v1 alongside a clear availability tradeoff. Across the 40 works (240 conditions) complete under both input apertures with the safety checker enabled, expanding the mask reduced PowerPaint v1 masked MAE by \textbf{44.7\%} (from 0.1055 to 0.0583; paired difference \ensuremath{-}0.0471 {[}\ensuremath{-}0.0589, \ensuremath{-}0.0360{]}), lowered local LPIPS from 0.2139 to 0.1714 (\ensuremath{-}0.0425 {[}\ensuremath{-}0.0593, \ensuremath{-}0.0248{]}), cut pixel-pair \texttt{\ensuremath{\Delta E_{00}}} from 11.0931 to 6.3145 (\ensuremath{-}4.7786 {[}\ensuremath{-}5.8002, \ensuremath{-}3.8261{]}), cut mean-Lab \texttt{\ensuremath{\Delta E_{00}}} from 8.5506 to 3.5543 (\ensuremath{-}4.9963 {[}\ensuremath{-}6.0864, \ensuremath{-}3.9571{]}), and reduced local DISTS from 0.1650 to 0.1430 (\ensuremath{-}0.0220 {[}\ensuremath{-}0.0340, \ensuremath{-}0.0096{]}). Gains concentrated where under-coverage was severest---\ensuremath{-}0.1494 on 1\% cracks and \ensuremath{-}0.1707 on 1\% flakes---while 5\% cracks (+0.0164) and 5\% contiguous losses (+0.0206) rose slightly; because grid expansion simultaneously alters explicit mask coverage and reduces surrounding visible context, these net changes reflect the combined input protocol rather than context reduction alone. Simultaneously, expanding the input aperture reduced PowerPaint v1\textquotesingle s scorable availability from \textbf{300/360 (83.3\%, 44/60 complete works)} to \textbf{281/360 (78.1\%, 41/60 complete works)}, with declines concentrated in oil on canvas (64 to 55/90) and oil on wood (70 to 64/90; 275 conditions passed under both apertures, 25 only with original masks, 6 only with expanded masks, and 54 blocked in both).

For SDXL across all 60 works (360 paired conditions), grid alignment left overall MAE virtually unchanged (0.3007 to 0.3015; +0.0008 {[}\ensuremath{-}0.0081, 0.0101{]}; LPIPS +0.0035 {[}\ensuremath{-}0.0030, 0.0101{]}; DISTS +0.0074 {[}0.0041, 0.0107{]}), because condition-level shifts varied in direction and canceled out: MAE decreased on 1\% cracks (\ensuremath{-}0.0704) and 1\% flakes (\ensuremath{-}0.0652), showed near-zero change on 5\% cracks (\ensuremath{-}0.0034, interval spanning zero), and increased on 1\% contiguous losses (+0.0269), 5\% flakes (+0.0594), and 5\% contiguous losses (+0.0574). On the exact 40-work (240-condition) intersection complete across all four models under both input apertures, mean masked MAE for LaMa, MAT, PowerPaint v1, and SDXL was \textbf{0.0274, 0.0284, 0.1055, and 0.3087} with original masks versus \textbf{0.0313, 0.0337, 0.0583, and 0.3115} with shared expanded masks, confirming that grid alignment substantially alters error magnitude without reversing LaMa \(\approx\) MAT \(\!<\!\) PowerPaint v1 \(\!<\!\) SDXL.

Our earlier pilot runs corroborate these full-cohort findings. On the 16-work native-scale pilot (seed 2026), nearest-neighbor downsampling missed 53\%--56\% of 1\% crack and flake pixels, and grid alignment reduced PowerPaint v1 MAE by 43.5\% (0.0974 to 0.0550; \ensuremath{-}0.0424 {[}\ensuremath{-}0.0617, \ensuremath{-}0.0251{]} on 12 complete works; \ensuremath{-}0.1252 on 1\% cracks and \ensuremath{-}0.1553 on 1\% flakes) and SDXL MAE from 0.2963 to 0.2730 (\ensuremath{-}0.0233 {[}\ensuremath{-}0.0374, \ensuremath{-}0.0104{]} on 16 works). Conversely, on the 60-work center-masked pilot---where the compact 5\% hole had only 7.19\% initial under-coverage---mask expansion slightly increased MAE for both SDXL (+0.0179 {[}0.0076, 0.0310{]}) and LaMa (+0.0072 {[}0.0043, 0.0107{]}) under the combined expanded input protocol.

\begin{figure}[!htb]
\centering
\includegraphics[width=\linewidth]{figures/figure-4-latent-coverage.pdf}
\caption{Latent-space mask coverage across one deterministic 1\% example each of a crack, scattered flakes, and a contiguous loss. Overlays distinguish original scoring-mask pixels reached versus missed (brick red) under 8\ensuremath{\times} nearest-neighbor downsampling, and the full coverage achieved after 8 \ensuremath{\times} 8 block expansion (green).}\label{figure:4}
\end{figure}

\subsection{\texorpdfstring{Prefill initialization strongly alters SDXL error and luminance when \texttt{strength\ \textless{}\ 1.0}}{Prefill initialization strongly alters SDXL error and luminance when strength \textless{} 1.0}}\label{prefill-initialization-strongly-alters-sdxl-error-and-luminance-when-strength-1.0}

Fixing SDXL at the model card\textquotesingle s example settings (\texttt{strength\ =\ 0.99}, CFG 8.0, 20 requested steps; 19 effective steps in Diffusers v0.35.2, 512 resolution, seed 2026) and varying \emph{only} the RGB color placed inside the expanded hole produced large, systematic error and brightness shifts in both protocols (Table~\ref{table:4}). On the 16-work native-scale pilot (96 conditions per color), replacing black prefill with an 8-pixel surrounding ring mean reduced masked MAE by more than one third, from 0.1513 to 0.0981 (paired difference \ensuremath{-}0.0532 {[}\ensuremath{-}0.1029, \ensuremath{-}0.0123{]}), lowered local LPIPS from 0.4208 to 0.3781 (\ensuremath{-}0.0427 {[}\ensuremath{-}0.0857, \ensuremath{-}0.0076{]}), and reversed signed luminance bias from \ensuremath{-}0.0847 (darkening) to +0.0339 (+0.1186 {[}0.0768, 0.1685{]}). Mid-gray prefill (RGB 128) gave intermediate MAE (0.1302) and positive luminance bias (+0.0553), trailing ring-color prefill by +0.0321 {[}0.0136, 0.0543{]} MAE. The univariate prefill effect peaked on 5\% contiguous losses, where ring-color prefill cut MAE from 0.2332 (black) to 0.1195 (\ensuremath{-}0.1137 {[}\ensuremath{-}0.2137, \ensuremath{-}0.0272{]}). Across all 60 center-masked works, the same pattern held: MAE fell from 0.2342 (black) to 0.1777 (mid-gray; \ensuremath{-}0.0566 {[}\ensuremath{-}0.0851, \ensuremath{-}0.0287{]}) and 0.1618 (ring mean; \ensuremath{-}0.0725 {[}\ensuremath{-}0.0998, \ensuremath{-}0.0475{]}), while luminance bias shifted from \ensuremath{-}0.1685 to +0.0251 and +0.0188. Across these two crop protocols, the overall black-to-ring MAE shifts alone (0.0532 and 0.0725) were five to seven times our 0.01 practical margin, with the 16-work condition mean for 5\% contiguous losses reaching 0.1137.

\begin{table}[!htb]
\caption{Univariate effect of hole prefill color on SDXL Inpainting (Protocol Factor 2). All arms fix \texttt{strength\ =\ 0.99}, CFG 8.0, 20 requested steps (19 effective), 512 resolution, and seed 2026. Signed luminance bias is output minus reference in normalized \texttt{{[}0,\ 1{]}} Rec. 709 sRGB (0 denotes neutral); down arrows (\ensuremath{\downarrow}) indicate lower error is better, and em dashes (---) mark the black-prefill reference rows.}\label{table:4}
\centering
\fontsize{10}{12}\selectfont
\setlength{\tabcolsep}{3pt}
\renewcommand{\arraystretch}{1.2}
\begin{tabular}{>{\raggedright\arraybackslash}p{\dimexpr0.180\linewidth-2.160\tabcolsep\relax}>{\raggedright\arraybackslash}p{\dimexpr0.115\linewidth-1.380\tabcolsep\relax}>{\raggedright\arraybackslash}p{\dimexpr0.105\linewidth-1.260\tabcolsep\relax}>{\raggedright\arraybackslash}p{\dimexpr0.105\linewidth-1.260\tabcolsep\relax}>{\raggedright\arraybackslash}p{\dimexpr0.160\linewidth-1.920\tabcolsep\relax}>{\raggedright\arraybackslash}p{\dimexpr0.335\linewidth-4.020\tabcolsep\relax}}
\toprule
\textbf{Protocol (works \ensuremath{\times} conditions)} & \textbf{Prefill color} & \textbf{Masked MAE \ensuremath{\downarrow}} & \textbf{Local LPIPS \ensuremath{\downarrow}} & \textbf{Signed luminance bias (0 neutral)} & \textbf{Paired MAE vs. black [95\% CI]} \\
\midrule
Center 5\% mask (60 \ensuremath{\times} 1) & Black (RGB 0) & 0.2342 & 0.4734 & \ensuremath{-}0.1685 & — \\
Center 5\% mask (60 \ensuremath{\times} 1) & Mid-gray (RGB 128) & 0.1777 & 0.4569 & +0.0251 & \ensuremath{-}0.0566 [\ensuremath{-}0.0851, \ensuremath{-}0.0287] \\
Center 5\% mask (60 \ensuremath{\times} 1) & 8-px ring mean & 0.1618 & 0.4566 & +0.0188 & \ensuremath{-}0.0725 [\ensuremath{-}0.0998, \ensuremath{-}0.0475] \\
Native-scale pilot (16 \ensuremath{\times} 6) & Black (RGB 0) & 0.1513 & 0.4208 & \ensuremath{-}0.0847 & — \\
Native-scale pilot (16 \ensuremath{\times} 6) & Mid-gray (RGB 128) & 0.1302 & 0.4059 & +0.0553 & \ensuremath{-}0.0211 [\ensuremath{-}0.0777, 0.0315] \\
Native-scale pilot (16 \ensuremath{\times} 6) & 8-px ring mean & 0.0981 & 0.3781 & +0.0339 & \ensuremath{-}0.0532 [\ensuremath{-}0.1029, \ensuremath{-}0.0123] \\
\bottomrule
\end{tabular}
\end{table}


Scaling this strongest SDXL configuration (\texttt{0.99\ /\ 8.0\ /\ 20} + ring-color prefill) to all 60 Cleveland works (360 conditions) substantially narrowed---but did not close---SDXL\textquotesingle s gap to non-diffusion baselines. Relative to SDXL\textquotesingle s main \texttt{1.0\ /\ 7.5\ /\ 25} configuration across all 60 works, ring-color \texttt{0.99} cut masked MAE by \ensuremath{-}0.1901 {[}\ensuremath{-}0.2354, \ensuremath{-}0.1460{]} (0.3015 to 0.1114), local LPIPS by \ensuremath{-}0.1104 {[}\ensuremath{-}0.1348, \ensuremath{-}0.0856{]} (0.4920 to 0.3816), pixel-pair \texttt{\ensuremath{\Delta E_{00}}} by \ensuremath{-}18.4006 {[}\ensuremath{-}22.6573, \ensuremath{-}14.2282{]} (31.3677 to 12.9671), mean-Lab \texttt{\ensuremath{\Delta E_{00}}} by \ensuremath{-}21.4536 {[}\ensuremath{-}25.6475, \ensuremath{-}17.2692{]} (30.0251 to 8.5715), and local DISTS by \ensuremath{-}0.0229 {[}\ensuremath{-}0.0336, \ensuremath{-}0.0119{]} (0.3438 to 0.3209). Even under this favorable setup, SDXL remained +0.0789 {[}0.0714, 0.0870{]} higher in MAE than LaMa and +0.0767 {[}0.0691, 0.0850{]} higher than MAT (LPIPS gaps +0.3153 {[}0.2923, 0.3397{]} and +0.3052 {[}0.2812, 0.3302{]}), with lower 95\% confidence bounds exceeding 0.01 in every medium-and-support group (+0.0763 oil on canvas, +0.0790 oil on wood, +0.0834 watercolor on paper, +0.0771 ink on paper).

Additional SDXL sensitivities on the 60 center-masked works confirmed that neither 1024-pixel inference nor removing the text prompt reversed the ranking against LaMa. Under \texttt{strength\ =\ 1.0}, upsampling to 1024 inference improved local LPIPS (\ensuremath{-}0.0728 {[}\ensuremath{-}0.0858, \ensuremath{-}0.0602{]}) and DISTS (\ensuremath{-}0.0503 {[}\ensuremath{-}0.0614, \ensuremath{-}0.0395{]}) while MAE uncertainty spanned zero (0.2647 to 0.2492; \ensuremath{-}0.0156 {[}\ensuremath{-}0.0342, 0.0029{]}), and an empty prompt reduced MAE by \ensuremath{-}0.0327 {[}\ensuremath{-}0.0522, \ensuremath{-}0.0133{]} (to 0.2320); both stayed well above shared-input LaMa (+0.1728 {[}0.1387, 0.2097{]} and +0.1519 {[}0.1261, 0.1801{]}). Pairing 1024 inference with \texttt{strength\ =\ 0.99} and black prefill \emph{worsened} MAE from 0.2492 to 0.3465 (+0.0974 {[}0.0300, 0.1603{]}) and intensified darkening (luminance bias \ensuremath{-}0.3483 vs. +0.0697 at \texttt{strength\ =\ 1.0}), even while 1024 versus 512 within the black-prefill model-card setup still lowered DISTS (\ensuremath{-}0.0284 {[}\ensuremath{-}0.0419, \ensuremath{-}0.0141{]})---showing how prefill sensitivity and inference resolution can pull pixel and perceptual metrics in opposite directions. Visual auditing of 12 randomly sampled Cleveland SDXL main outputs documented bright streaks across dark passages, smooth low-texture patches over granular paper or canvas, disrupted color transitions, and straight or triangular geometric bands absent from the originals, while two sampled cases reproduced pixel-for-pixel.

\subsection{Safety-checker withholding varies sharply across collections}\label{safety-checker-withholding-varies-sharply-across-collections}

Although PowerPaint v1 ran with identical weights, prompts, and safety-checker settings in both museums, its output availability differed sharply between collections and across media (Table~\ref{table:5}; Fig.~\ref{figure:6}). In Cleveland, PowerPaint returned scorable outputs for 281 of 360 planned conditions (78.1\%), with 79 conditions (21.9\%) across 19 works persistently withheld across both identical-seed attempts. Withholding was heaviest in oil on canvas (35/90 blocked; 61.1\% availability; 8/15 complete works) and oil on wood (26/90 blocked; 71.1\%; 8/15 complete), compared with watercolor on paper (78/90; 86.7\%; 12/15 complete) and ink on paper (84/90; 93.3\%; 13/15 complete). In Chicago, by contrast, PowerPaint achieved 97.2\% overall availability (350/360 conditions; 54/60 complete works), including 96.7\% in oil on canvas (87/90) and 100\% in watercolor on paper (90/90). All other models completed 360/360 conditions in both museums.

\begin{table}[!htb]
\caption{PowerPaint v1 output availability and conditional replay across museum strata (Protocol Factor 3). Results compare safety-checker-enabled runs across medium-and-support strata in Cleveland and Chicago with the single-pass safety-off replay of Cleveland\textquotesingle s 79 withheld conditions. Down arrows (\ensuremath{\downarrow}) indicate lower error is better, em dashes (---) mark stratum rows where only paired differences relative to LaMa are reported, and the asterisk (\texttt{*}) marks the partial-condition replay denominator explained in the table footnote.}\label{table:5}
\centering
\fontsize{10}{12}\selectfont
\setlength{\tabcolsep}{3pt}
\renewcommand{\arraystretch}{1.2}
\begin{tabular}{>{\raggedright\arraybackslash}p{\dimexpr0.220\linewidth-2.640\tabcolsep\relax}>{\raggedright\arraybackslash}p{\dimexpr0.120\linewidth-1.440\tabcolsep\relax}>{\raggedright\arraybackslash}p{\dimexpr0.120\linewidth-1.440\tabcolsep\relax}>{\raggedright\arraybackslash}p{\dimexpr0.120\linewidth-1.440\tabcolsep\relax}>{\raggedright\arraybackslash}p{\dimexpr0.120\linewidth-1.440\tabcolsep\relax}>{\raggedright\arraybackslash}p{\dimexpr0.300\linewidth-3.600\tabcolsep\relax}}
\toprule
\textbf{Collection \& stratum / arm} & \textbf{Scorable / planned conditions} & \textbf{Availability (\%)} & \textbf{Complete 6-condition works} & \textbf{PowerPaint MAE \ensuremath{\downarrow}} & \textbf{PowerPaint \ensuremath{-} LaMa MAE [95\% CI]} \\
\midrule
\textbf{Cleveland (safety checker on)} & \textbf{281 / 360} & \textbf{78.1\%} & \textbf{41 / 60} & \textbf{0.0581} & \textbf{+0.0266 [0.0203, 0.0337]} \\
— Oil on canvas & 55 / 90 & 61.1\% & 8 / 15 & — & +0.0394 [0.0251, 0.0542] \\
— Oil on wood & 64 / 90 & 71.1\% & 8 / 15 & — & +0.0336 [0.0250, 0.0422] \\
— Watercolor on paper & 78 / 90 & 86.7\% & 12 / 15 & — & +0.0125 [0.0086, 0.0167] \\
— Ink on paper & 84 / 90 & 93.3\% & 13 / 15 & — & +0.0274 [0.0142, 0.0437] \\
\textbf{Cleveland conditional replay (safety off)} & \textbf{79 / 79} & \textbf{100.0\%} & \textbf{—*} & \textbf{0.0714} & \textbf{+0.0369 [0.0190, 0.0583]} \\
\textbf{Cleveland combined mixed cohort} & \textbf{360 / 360} & \textbf{100.0\%} & \textbf{60 / 60} & \textbf{0.0604} & \textbf{+0.0279 [0.0211, 0.0357]} \\
\textbf{Chicago (safety checker on)} & \textbf{350 / 360} & \textbf{97.2\%} & \textbf{54 / 60} & \textbf{0.0650} & \textbf{+0.0263 [0.0221, 0.0310]} \\
— Oil on canvas & 87 / 90 & 96.7\% & 13 / 15 & — & +0.0215 [0.0165, 0.0274] \\
— Oil on wood & 86 / 90 & 95.6\% & 12 / 15 & — & +0.0318 [0.0238, 0.0409] \\
— Watercolor on paper & 90 / 90 & 100.0\% & 15 / 15 & — & +0.0312 [0.0220, 0.0422] \\
— Ink on paper & 87 / 90 & 96.7\% & 14 / 15 & — & +0.0210 [0.0134, 0.0295] \\
\bottomrule
\end{tabular}
\par\smallskip
\begin{minipage}{\linewidth}\fontsize{10}{12}\selectfont
* Each of the 19 affected Cleveland works replayed only its own withheld conditions (between 1 and 6 conditions per work; all 6 conditions were withheld and replayed for \emph{The Trinity}), totaling 79 replayed conditions. Replay summary scores average first within each work across its replayed conditions and then weight all 19 works equally.
\end{minipage}
\end{table}


Museum catalog descriptions for the 19 affected Cleveland paintings showed that 12 centered on human figures (portraits and narrative scenes; 58 withheld conditions), 3 on animals (6 conditions), and 4 on landscapes, woodlands, or flora (15 conditions)---referring to full-artwork catalog records, including \emph{The Trinity}, whose uniformly black crop was blocked on all six conditions. Because these counts reflect whole-artwork museum descriptions whereas the safety checker evaluates individual generated crops, catalog subjects cannot identify why the checker flagged any particular output.

Replaying the exact 79 withheld Cleveland conditions once with the safety checker disabled recovered 79/79 outputs. Because each of the 19 affected works replayed only its own withheld conditions (from 1 up to all 6 for \emph{The Trinity}), we averaged first within each work across its replayed conditions and then weighted all 19 affected works equally. Under this work-level aggregation, PowerPaint yielded a masked MAE of 0.0714 and a paired difference relative to LaMa of +0.0369 {[}0.0190, 0.0583{]} (+0.1263 {[}0.0714, 0.1898{]} local LPIPS), compared with 0.0581 and +0.0266 {[}0.0203, 0.0337{]} across the 41 works complete over all six safety-on conditions (246 conditions); each estimate reflects its own work and condition denominator. Combining the 281 original passed outputs with the 79 replayed outputs across all 60 works gave a PowerPaint MAE of 0.0604 (+0.0279 {[}0.0211, 0.0357{]} vs. LaMa 0.0325; +0.1032 {[}0.0820, 0.1257{]} LPIPS), alongside MAT at 0.0347 and SDXL main at 0.3015. Recovering the missing outputs slightly increased PowerPaint\textquotesingle s overall MAE while preserving LaMa \(\approx\) MAT \(\!<\!\) PowerPaint v1 \(\!<\!\) SDXL.

\begin{figure}[!htb]
\centering
\includegraphics[width=\linewidth]{figures/figure-6-powerpaint-coverage.pdf}
\caption{PowerPaint v1 scorable output availability across medium-and-support groups in both museums. Horizontal stacked bars show scorable versus safety-withheld conditions (brick red) out of 90 planned conditions per group with the built-in safety checker enabled, yielding 78.1\% overall availability (281/360) in the Cleveland Museum of Art and 97.2\% (350/360) in the Art Institute of Chicago.}\label{figure:6}
\end{figure}

\subsection{Rankings held across damage geometries and painting media, whereas between-group contrasts varied across collections}\label{rankings-held-across-damage-geometries-and-painting-media-whereas-between-group-contrasts-varied-across-collections}

Stratifying results by procedural damage morphology and area fraction confirms identical model ordering across all 12 condition-by-museum cells (Table~\ref{table:6}; Fig.~\ref{figure:5}). In every combination of damage type (crack, scattered flakes, contiguous loss), area fraction (1\%, 5\%), and museum (Cleveland, Chicago), masked MAE followed \textbf{LaMa \(\approx\) MAT \(\!<\!\) PowerPaint v1 \(\!<\!\) SDXL}. For LaMa, MAT, and PowerPaint v1, mean error was consistently higher at 5\% than at 1\% within each damage family, peaking on 5\% contiguous losses in both collections (0.0345, 0.0390, and 0.0697 in Cleveland; 0.0477, 0.0558, and 0.0956 in Chicago). SDXL under \texttt{strength\ =\ 1.0} behaved differently, producing persistently high MAE across all geometries (0.2874--0.3359 in Cleveland; 0.2670--0.3003 in Chicago) and peaking on 1\% cracks (0.3359 and 0.3003) despite their small area.

\begin{table}[!htb]
\caption{Masked MAE stratified by procedural damage type and area fraction. Values are computed under the shared 8 \ensuremath{\times} 8 expanded input aperture and averaged over all artworks complete across all four models for each condition; \texttt{n} denotes the number of matched works per condition, and down arrows (\ensuremath{\downarrow}) indicate lower error is better.}\label{table:6}
\centering
\fontsize{10}{12}\selectfont
\setlength{\tabcolsep}{3pt}
\renewcommand{\arraystretch}{1.2}
\begin{tabular}{>{\raggedright\arraybackslash}p{\dimexpr0.110\linewidth-1.540\tabcolsep\relax}>{\raggedright\arraybackslash}p{\dimexpr0.180\linewidth-2.520\tabcolsep\relax}>{\raggedright\arraybackslash}p{\dimexpr0.120\linewidth-1.680\tabcolsep\relax}>{\raggedright\arraybackslash}p{\dimexpr0.135\linewidth-1.890\tabcolsep\relax}>{\raggedright\arraybackslash}p{\dimexpr0.125\linewidth-1.750\tabcolsep\relax}>{\raggedright\arraybackslash}p{\dimexpr0.155\linewidth-2.170\tabcolsep\relax}>{\raggedright\arraybackslash}p{\dimexpr0.175\linewidth-2.450\tabcolsep\relax}}
\toprule
\textbf{Collection} & \textbf{Damage condition} & \textbf{Matched works (n)} & \textbf{LaMa MAE \ensuremath{\downarrow}} & \textbf{MAT MAE \ensuremath{\downarrow}} & \textbf{PowerPaint v1 MAE \ensuremath{\downarrow}} & \textbf{SDXL (main) MAE \ensuremath{\downarrow}} \\
\midrule
Cleveland & Crack 1\% & 45 & 0.0290 & 0.0299 & 0.0535 & 0.3359 \\
Cleveland & Crack 5\% & 47 & 0.0312 & 0.0339 & 0.0641 & 0.2961 \\
Cleveland & Flake 1\% & 47 & 0.0296 & 0.0302 & 0.0488 & 0.3004 \\
Cleveland & Flake 5\% & 45 & 0.0309 & 0.0331 & 0.0568 & 0.3042 \\
Cleveland & Contiguous loss 1\% & 49 & 0.0341 & 0.0366 & 0.0470 & 0.2874 \\
Cleveland & Contiguous loss 5\% & 48 & 0.0345 & 0.0390 & 0.0697 & 0.2887 \\
Chicago & Crack 1\% & 58 & 0.0302 & 0.0314 & 0.0514 & 0.3003 \\
Chicago & Crack 5\% & 58 & 0.0377 & 0.0417 & 0.0711 & 0.2670 \\
Chicago & Flake 1\% & 58 & 0.0312 & 0.0315 & 0.0500 & 0.2864 \\
Chicago & Flake 5\% & 58 & 0.0362 & 0.0391 & 0.0620 & 0.2823 \\
Chicago & Contiguous loss 1\% & 58 & 0.0404 & 0.0431 & 0.0628 & 0.2694 \\
Chicago & Contiguous loss 5\% & 60 & 0.0477 & 0.0558 & 0.0956 & 0.2954 \\
\bottomrule
\end{tabular}
\end{table}


\begin{figure}[!htb]
\centering
\includegraphics[width=\linewidth]{figures/figure-5-condition-mae.pdf}
\caption{Masked MAE across the six procedural damage conditions in Cleveland and Chicago. Points show complete-case means with 10,000-resample artwork-level bootstrap 95\% percentile intervals on a logarithmic vertical axis for LaMa, MAT, PowerPaint v1, and SDXL Inpainting under the shared 8 \ensuremath{\times} 8 expanded aperture; axis annotations give the number of artworks with scorable outputs from all four models per condition.}\label{figure:5}
\end{figure}

Stratifying by physical painting medium and support likewise preserved LaMa \(\approx\) MAT \(\!<\!\) PowerPaint v1 \(\!<\!\) SDXL across all eight museum-by-medium strata (four-model common complete sample sizes: 8, 8, 12, and 13 works in Cleveland; 13, 12, 15, and 14 in Chicago). In model-specific paired comparisons against LaMa (Table~\ref{table:7}), MAT remained within our 0.01 practical margin in every group (+0.0013 to +0.0045 MAE; every upper 95\% confidence bound \ensuremath{\leq} 0.0071), whereas PowerPaint v1 (+0.0125 to +0.0394) and SDXL (+0.1990 to +0.3447) had strictly positive paired MAE differences with confidence intervals above zero.

\begin{table}[!htb]
\caption{Artwork-level paired masked MAE differences relative to LaMa stratified by painting medium and support. Entries show candidate minus LaMa differences with 95\% bootstrap percentile CIs (positive values indicate higher error than LaMa). MAT and SDXL comparisons use all 15 works per group; PowerPaint v1 uses complete six-condition works (\texttt{n} = 8, 8, 12, and 13 in Cleveland; \texttt{n} = 13, 12, 15, and 14 in Chicago).}\label{table:7}
\centering
\fontsize{10}{12}\selectfont
\setlength{\tabcolsep}{3pt}
\renewcommand{\arraystretch}{1.2}
\begin{tabular}{>{\raggedright\arraybackslash}p{\dimexpr0.110\linewidth-1.100\tabcolsep\relax}>{\raggedright\arraybackslash}p{\dimexpr0.150\linewidth-1.500\tabcolsep\relax}>{\raggedright\arraybackslash}p{\dimexpr0.245\linewidth-2.450\tabcolsep\relax}>{\raggedright\arraybackslash}p{\dimexpr0.250\linewidth-2.500\tabcolsep\relax}>{\raggedright\arraybackslash}p{\dimexpr0.245\linewidth-2.450\tabcolsep\relax}}
\toprule
\textbf{Collection} & \textbf{Medium and support} & \textbf{MAT \ensuremath{-} LaMa MAE [95\% CI]} & \textbf{PowerPaint v1 \ensuremath{-} LaMa MAE [95\% CI]} & \textbf{SDXL (main) \ensuremath{-} LaMa MAE [95\% CI]} \\
\midrule
Cleveland & Oil on canvas & +0.0013 [0.0006, 0.0021] & +0.0394 [0.0251, 0.0542] & +0.3378 [0.2346, 0.4416] \\
Cleveland & Oil on wood & +0.0021 [0.0010, 0.0035] & +0.0336 [0.0250, 0.0422] & +0.3317 [0.2335, 0.4407] \\
Cleveland & Watercolor on paper & +0.0024 [0.0016, 0.0033] & +0.0125 [0.0086, 0.0167] & +0.2076 [0.1352, 0.2820] \\
Cleveland & Ink on paper & +0.0030 [0.0011, 0.0057] & +0.0274 [0.0142, 0.0437] & +0.1990 [0.1155, 0.2904] \\
Chicago & Oil on canvas & +0.0023 [0.0012, 0.0035] & +0.0215 [0.0165, 0.0274] & +0.2286 [0.1638, 0.3032] \\
Chicago & Oil on wood & +0.0035 [0.0018, 0.0058] & +0.0318 [0.0238, 0.0409] & +0.3447 [0.2554, 0.4367] \\
Chicago & Watercolor on paper & +0.0022 [0.0010, 0.0036] & +0.0312 [0.0220, 0.0422] & +0.2059 [0.1290, 0.2944] \\
Chicago & Ink on paper & +0.0045 [0.0023, 0.0071] & +0.0210 [0.0134, 0.0295] & +0.2129 [0.1185, 0.3211] \\
\bottomrule
\end{tabular}
\end{table}


In contrast to this ranking stability, the \emph{magnitude} of between-medium differences did not replicate stably across collections. In Cleveland, SDXL\textquotesingle s MAE gap relative to LaMa appeared smaller on watercolor on paper than on oil on canvas by \ensuremath{-}0.1303 {[}\ensuremath{-}0.2562, \ensuremath{-}0.0039{]}. In Chicago, that same contrast shrank to \ensuremath{-}0.0228 {[}\ensuremath{-}0.1303, 0.0877{]}, spanning zero, while oil on wood showed the largest SDXL gap (+0.3447 vs. +0.2286 on oil on canvas). Similarly, PowerPaint\textquotesingle s smallest gap relative to LaMa occurred in watercolor on paper in Cleveland (+0.0125) but in ink on paper (+0.0210) and oil on canvas (+0.0215) in Chicago. Model ordering therefore held across the tested strata, whereas cross-group gap magnitudes remained collection-specific.

\section{Discussion}\label{discussion}

\subsection{Guidance for museum and digital-heritage practitioners}\label{guidance-for-museum-and-digital-heritage-practitioners}

Our two-museum results provide a clear operational rule for practitioners selecting off-the-shelf inpainting tools for localized painting losses (1\%--5\% area): \textbf{model choice did not depend on painting medium or support.} Across oil on canvas, oil on wood, watercolor on paper, and ink on paper in both Cleveland and Chicago, LaMa and MAT delivered practically interchangeable full-reference fidelity (differing by at most 0.0045 MAE and 0.0123 LPIPS in stratum point estimates, with overall differences inside our 0.01 MAE and 0.02 LPIPS margins), required no prompt engineering, and achieved 100\% output availability without safety-filter refusals. On fine cracks, dispersed flakes, and small lacunae, these pixel-space architectures consistently tracked local palette and surrounding texture more faithfully than the tested latent diffusion models.

When latent diffusion pipelines are considered---for example, to explore generative hypotheses over wider semantic losses---our experiments highlight two practical deployment hazards. First, default pipeline settings can severely distort local color: SDXL at \texttt{strength\ =\ 1.0} produced mean-Lab \texttt{\ensuremath{\Delta E_{00}}} shifts above 26, and \texttt{strength\ =\ 0.99} with black prefill introduced substantial darkening (\ensuremath{-}0.0847 to \ensuremath{-}0.1685 luminance). Even under its strongest tested ring-color configuration across Cleveland, SDXL\textquotesingle s masked MAE (0.1114) remained more than three times that of LaMa (0.0325). Second, built-in safety checkers can silently block large shares of museum imagery---nearly 39\% of oil-on-canvas conditions in Cleveland---turning batch workflows into unpredictable subset processors unless refusal rates are monitored explicitly.

\subsection{Reporting standards for painting-inpainting benchmarks}\label{reporting-standards-for-painting-inpainting-benchmarks}

For researchers designing or reviewing digital restoration benchmarks, our paired comparisons show that mask alignment and SDXL initialization settings can shift masked MAE by 0.047 to 0.190---far above our 0.01 practical margin---even with identical model weights and test artworks, while safety filtering alters both the scored sample and availability (including when mask expansion itself reduces PowerPaint availability from 83.3\% to 78.1\%). We therefore recommend three mandatory reporting items for any full-reference painting-inpainting benchmark:

\begin{enumerate}
\def\labelenumi{\arabic{enumi}.}
\tightlist
\item
  \textbf{Input-to-latent mask alignment and compositing boundary.} Specify how pixel masks map onto the latent grid, report the fraction of target pixels covered (or the expanded-to-original area ratio if aligned), composite generated pixels strictly inside the original unexpanded mask, and include an original-mask sensitivity check for pixel-space baselines so that aperture tradeoffs remain transparent.
\item
  \textbf{Pre-masking and denoising initialization.} Whenever \texttt{strength\ \textless{}\ 1.0}, verify that ground-truth pixels are erased prior to VAE encoding and state the exact prefill rule (black, mid-gray, or surrounding-ring mean). Reporting \texttt{strength\ =\ 1.0} alongside explicit prefill-color sensitivities prevents arbitrary initialization choices from being conflated with model capability.
\item
  \textbf{Planned denominators and refusal accounting.} Report planned, scorable, and withheld counts per model and stratum alongside complete-case fidelity scores. When safety checkers or runtime failures exclude a non-trivial share of conditions, include conditional sensitivity checks to verify whether selective attrition shifts error estimates or rankings.
\end{enumerate}

Committing pre-specified sampling orders, frozen crop and mask hashes, code-provenance metadata, and practical decision margins to version control before inference further guards against inadvertent selection bias. As our Cleveland--Chicago contrast illustrates, an apparent between-medium difference in a single collection (watercolor versus oil on canvas) can shrink and lose precision upon independent replication.

\subsection{Limitations of the study}\label{limitations-of-the-study}

Five scope boundaries govern our findings and point directly to future work. First, all damage masks were procedurally generated (cracks, scattered flakes, and contiguous losses at 1\% and 5\% area) rather than authentic historical loss contours transplanted from damaged donor paintings; authentic mask-transfer studies require sufficient medium- and support-matched donor annotations with clear provenance and reuse licenses, which remained scarce in our audit of ARTeFACT \citep{ivanova2024artefact}. In addition, our main runs used a single fixed inference seed (\texttt{2026}) and 512-pixel resolution across all four pinned models, whereas SDXL was trained at 1024 \ensuremath{\times} 1024; although our 1024 center-masked sensitivity did not alter the ranking against LaMa, multi-seed native-1024 studies and newer rectified-flow models such as FLUX.1 Fill \citep{fluxfill} on higher-memory or quantized setups warrant separate investigation. Second, because 512 \ensuremath{\times} 512 crops were taken without further resampling from each museum\textquotesingle s published rendition, Cleveland\textquotesingle s full-size files and Chicago\textquotesingle s 1686-pixel-wide IIIF JPEGs may span different physical painting areas per pixel; physical scale standardization requires per-artwork millimeter calibration from museum imaging metadata. Third, painting media and supports naturally covary with subject matter, historical period, brushwork scale, and digitization pipelines, so between-group contrasts are descriptive rather than causal effects of material alone. Fourth, 16 of the 60 Cleveland paintings overlapped with our earlier native-scale pilot, which we mitigated through the independently sampled Chicago replication cohort while retaining two frozen edge cases (\emph{The Trinity} black crop and Chicago\textquotesingle s card-mounted watercolor \texttt{aic-114337}) verified via exclusion checks. Fifth, automatic full-reference pixel, color, and perceptual distances measure agreement with a digitized photograph under synthetic masking, not historical authenticity or conservation validity; physical restoration decisions remain the responsibility of trained conservators guided by archival and technical evidence \citep{aic1994, icomcc2008}.

\section{Conclusions}\label{conclusions}

Across 120 open-access paintings from two independent museums, three evaluation protocol choices---nearest-neighbor latent mask downsampling, hole prefill color under \texttt{strength\ \textless{}\ 1.0}, and safety-filter denominator accounting---substantially alter measured reconstruction error and effective sample composition in AI painting inpainting. Under the controlled protocol, the four-model masked-MAE ordering (LaMa \(\approx\) MAT \(\!<\!\) PowerPaint v1 \(\!<\!\) SDXL, with LaMa and MAT within our pre-specified practical margin) was identical on matched cases across both museum collections, all four painting media and supports, and all six procedural damage conditions; these three choices changed error magnitudes and the scored sample but did not reverse the ordering in our tests. For localized losses, strong non-diffusion baselines provided a consistent choice across painting media, whereas future heritage-inpainting benchmarks must report mask alignment, prefill initialization, and coverage denominators to produce comparable, trustworthy evidence.

\FloatBarrier

\section*{Acknowledgments}

During the preparation of this manuscript/study, the authors used Claude (Anthropic), Gemini (Google), and Codex (OpenAI) \emph{(specific tool version details to be confirmed by the author)} for protocol review and figure generation, English drafting and editing, and experimental coding, figure generation, literature verification, and statistical verification, respectively. The authors take responsibility for the content of this publication. \emph{(Author to confirm tool version details, scope of personal review, and any additional acknowledgments prior to submission.)}

\section*{Additional Information and Declarations}

\subsection*{Funding}
This research received no external funding. \emph{(To be confirmed by the author prior to submission.)}

\subsection*{Competing Interests}
The authors declare no conflicts of interest. \emph{(To be confirmed by the author prior to submission.)}

\subsection*{Author Contributions}
All authors approved the manuscript and agree to its submission to \textit{PeerJ Computer Science}. \emph{(Individual contributor roles to be confirmed by the authors.)}

\subsection*{Data Availability}
Source painting images are publicly accessible open-access museum records: 60 CC0 works from the Cleveland Museum of Art Open Access Collection \citetext{\citeauthor{cma-openaccess}, n.d.} and 60 public-domain works obtained via the Art Institute of Chicago Open Access / IIIF Image API (1686-pixel-wide JPEG renditions) \citetext{\citeauthor{aic-openaccess}, n.d.}. The public reproduction repository includes the experimental protocols, museum acquisition and crop preparation, mask generation, model adapters, scoring, analysis, and figure code. \url{https://github.com/liangdagongjue-code/mengwei} \emph{(Project license and publication release terms to be confirmed by the authors.)} To respect repository size limits and third-party distribution hygiene, raw museum downloads, per-artwork manifests, frozen per-condition \texttt{prepared.json} files (recording complete object IDs, sampling order, crop coordinates, and per-mask hashes), full JSON score summaries, contact sheets, and per-condition output PNGs are maintained in a separate experiment server archive rather than redistributed in Git \emph{(Before submission, provide the archival DOI for the reproducible code and machine-readable data, or identify the publishable supplemental files and any access arrangements approved by the journal.)}. The experiment server archive also contains run logs and per-execution provenance sidecars where recorded; historical runs executed outside a Git checkout may contain null Git commit fields, as documented in \nolinkurl{docs/RUN_PROVENANCE.md}. Third-party model weights remain governed by their respective upstream release artifacts and terms: LaMa (IOPaint TorchScript release \texttt{big-lama.pt} \citep{suvorov2021, lamaweights}), MAT (IOPaint Places-512 state dictionary \nolinkurl{Places_512_FullData_G.pth} \citep{li2022mat, matweights}), PowerPaint v1 (pinned author source revision \texttt{5b4c3d5} and weight revision \texttt{b4174b7} \citep{zhuang2023, powerpaintv1src, powerpaintv1weights}), and SDXL Inpainting 0.1 (pinned weight revision \texttt{115134f} under the CreativeML Open RAIL++-M License, with intended research use documented on the model card \citep{sdxlinpaint, sdxlinpaintweights}).

\subsection*{Supplemental Information}
The supplemental PDF contains Supplementary Note S1: Museum Cohort Screening, Exclusions, and Mask Specifications; Supplementary Note S2: Pinned Model Checkpoints, Environments, and Pre-Screening; Supplementary Note S3: SDXL Dual-Path Initialization Scaling and Audits; and Table S1: 16-Work Native-Scale Pilot Latent-Alignment Summary.

\section*{Abbreviations}
\noindent\begin{tabular}{@{}lp{0.75\linewidth}@{}}
AI & Artificial intelligence \\
AIC & Art Institute of Chicago \\
CC0 & Creative Commons Zero \\
CFG & Classifier-free guidance \\
CI & Confidence interval \\
CIELAB & CIE 1976 L*a*b* color space \\
CMA & Cleveland Museum of Art \\
DDIM & Denoising Diffusion Implicit Models \\
DISTS & Deep Image Structure and Texture Similarity \\
EXIF & Exchangeable Image File Format \\
GPU & Graphics processing unit \\
ICOM-CC & International Council of Museums Committee for Conservation \\
IIIF & International Image Interoperability Framework \\
JPEG & Joint Photographic Experts Group \\
LaMa & Large Mask Inpainting \\
LPIPS & Learned Perceptual Image Patch Similarity \\
MAE & Mean absolute error \\
MAT & Mask-Aware Transformer \\
RGB & Red–green–blue \\
SDXL & Stable Diffusion XL \\
SHA-256 & Secure Hash Algorithm 256 \\
sRGB & Standard RGB \\
VAE & Variational autoencoder \\
VRAM & Video random-access memory \\
\ensuremath{\Delta E_{00}} & CIEDE2000 perceptual color difference
\end{tabular}

\bibliographystyle{apalike}

\bibliography{references}

\end{document}
